\section{The Probabilistic Audit: Certifying Pipelines at Scale}
\label{sec:audit}

The theoretical guarantees in Section~\ref{sec:theorems} and the learning equivalences developed in Section~\ref{sec:preference} rely on the assumption that the sufficiency, merge-consistency, and idempotence conditions hold. In a real deployment, we cannot simply assume these properties; we must verify them. However, checking every node in a massive tree defeats the purpose of summarization.

We solve this dilemma using a \emph{probabilistic audit}. Because Theorem~\ref{thm:one-pass} guarantees that global fidelity is structurally implied by local fidelity, we do not need to validate the final summary against the full text to monitor quality. Instead, we treat the summarization pipeline as a stochastic process and estimate the \emph{Empirical Violation Rate} of each condition. A fixed sampling budget per audit run suffices to estimate these \emph{rates}; certifying that a specific book is error-free still requires those rates to be extremely small.

Throughout we assume the atomic objects we audit---leaf blocks $b$ and stored child pairs $(s_{u_L},s_{u_R})$---are small enough that both the oracle $\fstar$ and the summarizer $g$ can process them directly. The hierarchy is useful precisely because the root document $X$ may be orders of magnitude larger than this audit budget; we audit the bounded pieces and rely on the structural results to propagate guarantees upward.

\subsection{The Audit Protocol}

The audit reports three sample means: the sufficiency violation rate ($\hat{p}_{\text{suff}}$), the merge-consistency violation rate ($\hat{p}_{\text{assoc}}$), and the idempotence violation rate ($\hat{p}_{\text{idem}}$).

\paragraph{Sampling Leaves (Sufficiency Condition).}
We sample $n_\ell$ leaf blocks $\{b_i\}$ uniformly from the corpus. For each block, we generate its summary $g(b_i)$ and compare the oracle values:
\[
\hat{p}_{\text{suff}} := \frac{1}{n_\ell} \sum_{i=1}^{n_\ell} \mathbb{I}\left\{ \metric\big(\fstar(g(b_i)), \fstar(b_i)\big) > \tau \right\},
\]
where $\tau$ is a tolerance threshold (usually 0). This statistic is the empirical violation rate of Condition~\eqref{eq:leaf_sufficiency}. The check is cheap because leaf blocks are by definition short enough to fit in the model context.

\paragraph{Sampling Merge Consistency (The Structural Check).}
Condition~\eqref{eq:leaf_compatibility} requires a chain of equivalences, but we rarely have enough context to verify the entire chain at every node. Instead, we audit whichever link fits in memory for the sampled pair, always conditioning on the realized structure.
\begin{itemize}
    \item \textbf{Case A: Leaf boundaries (Substitution).} When $u,v$ are adjacent raw blocks, $u\concat v$ fits in context. We compare the joint and disjoint summaries:
    \[
    I_{\text{bound}} := \mathbb{I}\Big\{ \metric\big(\fstar(g(u\concat v)),\,\fstar(g(g(u)\concat g(v)))\big) > \tau \Big\}.
    \]
    \item \textbf{Case B: Internal merges (Compression).} When $u,v$ are stored summaries, they are short. We compare the raw concatenation of the stored strings to the one-shot summary:
    \[
    I_{\text{merge}} := \mathbb{I}\Big\{ \metric\big(\fstar(u\concat v),\,\fstar(g(u\concat v))\big) > \tau \Big\}.
    \]
\end{itemize}
For a batch of sampled leaf boundaries $\mathcal{B}$ and internal merges $\mathcal{M}$, the merge-consistency violation rate is the weighted average
\[
\hat{p}_{\text{assoc}} := \frac{\sum_{(u,v)\in\mathcal{B}} I_{\text{bound}}(u,v) + \sum_{(u,v)\in\mathcal{M}} I_{\text{merge}}(u,v)}{|\mathcal{B}| + |\mathcal{M}|}.
\]

\paragraph{Sampling Stability (Idempotence Condition).}
If the pipeline includes post-processing rounds, we sample $n_s$ summaries $z_k$ that have already passed through $g$ and check if re-summarization alters the oracle:
\[
\hat{p}_{\text{idem}} := \frac{1}{n_s} \sum_{k=1}^{n_s} \mathbb{I}\left\{ \metric\big(\fstar(g(z_k)), \fstar(z_k)\big) > \tau \right\}.
\]

These three sample means provide unbiased estimates of the violation probabilities. Because every check fits in context, teams can run them on demand, streaming new samples until the observed rate and its confidence interval fall below a deployment threshold.

\subsection{Scaling with Tree Size}

Let a realized hierarchy contain $N$ leaves and $M=N-1$ merges. A transparent union bound (Appendix~\ref{app:practical_audit}) then gives
\begin{equation}\label{eq:error_budget}
\Pr\big[\text{root violation}\big]\ \le\ N\,p_{\text{suff}} + M\,p_{\text{assoc}} + (R-1)\,p_{\text{idem}},
\end{equation}
where $p_{\text{suff}},p_{\text{assoc}},p_{\text{idem}}$ denote the true per-edge violation probabilities and $R$ counts optional clean-up rounds. This $N\times p$ scaling is the price of safety: even tiny per-edge failure rates accumulate over thousands of nodes. The audit's job is therefore to drive each $p$ low enough that the product $N\times p$ (and $M\times p$) stays within tolerance, or to change the decomposition until it does.

It bears repeating that these statistics describe the \emph{edge failure rates}, not the correctness of a specific document. If a document has $N$ leaves and the sufficiency failure rate is $p_{\text{suff}}$, the probability that every leaf is correct is roughly $(1-p_{\text{suff}})^N \approx e^{-Np_{\text{suff}}}$. For large $N$, the root summary is reliable only when the local error rates are correspondingly tiny.

\subsection{Implementation in ThinkingTrees}

The \textsc{ThinkingTrees} package (Section~\ref{sec:software}) implements the probabilistic audit through the \texttt{TreeAuditor} class. The key features are:

\begin{itemize}[nosep]
    \item \textbf{Hoeffding bounds}: Confidence intervals for violation rates using the Hoeffding inequality, providing rigorous statistical guarantees.
    \item \textbf{Configurable sampling}: Users specify sample budgets for each condition type (leaves, merges, idempotence).
    \item \textbf{Streaming mode}: Samples can be drawn incrementally until confidence intervals tighten below tolerance.
    \item \textbf{Audit flags}: Nodes exceeding violation thresholds are flagged for human review.
\end{itemize}

Appendix~\ref{app:practical_audit} provides step-by-step audit procedures including sample-size calculations.

\subsection{Optimizing the Pipeline}

The audit doubles as supervision for prompt engineering. We instantiate two promptable modules (leaf and merge), treat their instructions as parameters $\theta$, and update them using DSPy \citep{khattab2024dspy}. This is a form of active learning: violations reveal exactly which examples to upweight next.

\begin{enumerate}
    \item \textbf{Collect leaf data.} Sample leaves, run the current leaf module, and evaluate Condition~\eqref{eq:leaf_sufficiency}. Successful generations $(b, s_b)$ are appended to the module's context. When a failure occurs, we may re-generate at higher diversity or request a human edit.
    \item \textbf{Collect merge data.} Using the best available leaf summaries, sample adjacent pairs, run the merge module, and evaluate Condition~\eqref{eq:leaf_compatibility}. Successful merges become demonstrations for future runs.
    \item \textbf{Update and stop.} Re-optimize the prompt parameters via DSPy's gradient-free controller on the accumulated traces and resume sampling. Continue until the estimated violation rate falls below tolerance.
\end{enumerate}

This ``consistency bootstrap'' requires no external supervision beyond the oracle already assumed by the audit. Appendix~\ref{app:bootstrap} expands each step into DSPy-ready pseudocode.

\paragraph{Calibration vs.\ deployment.}
The audit statistics assume the summarizer $g$ is fixed. In practice we therefore separate the \emph{calibration phase}, during which prompts or weights are updated using the bootstrap loop, from the \emph{deployment phase}, during which we freeze $g$ and estimate $(\hat{p}_{\text{suff}},\hat{p}_{\text{assoc}},\hat{p}_{\text{idem}})$ on held-out data.
